%!TEX root = ./../../thesis-trier.tex

Here, in \cref{constr_sync:thm:finite}, 
we prove that for finite $L \subseteq \Sigma^*$, which are obviously sparse, the problem \constrsync{L} is in $\AC^0$.
In \cref{constr_sync:prop:finite} we prove that removing finite sets
from a constraint language does not lower the complexity if the resulting language is non-empty.

\begin{theorem}
\label{constr_sync:thm:finite}
	Let $L \subseteq \Sigma^*$ be finite.
	Then \constrsync{L} is in $\AC^0$
	(and in constant time for $L = \emptyset$ or $L = \{\varepsilon\}$). 
\end{theorem}
\begin{proof}
	An algorithm for \constrsync{\emptyset} always
	returns {\sc NO}, and for \constrsync{\{\varepsilon\}}
	it returns {\sc YES} if and only if the input
	automaton has a single state and {\sc NO} otherwise;
	both algorithms run in constant time.
	Next, let $u \in \Sigma^*$ and consider
	the morphism $\varphi \colon \{a\}^* \to \Sigma^*$ 
	given by $\varphi(a) = u$.
	By \cref{thm:hom_lower_bound_complexity},
	we get $\constrsync{\{u\}} \leACZEROm \constrsync{\{a\}}$,
	where the latter problem is in $\AC^0$ by \crefitem{constr_sync:thm:introductory_examples}{constrsync:ex:item:a}.
	So, $\constrsync{\{u\}}$ for $u \in \Sigma^*$
	is in $\AC^0$.
	Then, for finite~$L$, the claim follows from 
	\cref{constr_sync:thm:union}.	
\end{proof} 

More precisely, by closer inspection of the proof of \cref{constr_sync:thm:finite}
(and the proofs of the referenced statements), we obtain that finite sets can be handled
inside of $\leACZEROm$-reductions, giving the following statement.

\begin{proposition} \label{constr_sync:prop:finite}
	Let $L \subseteq \Sigma^*$ be regular and $F \subseteq \Sigma^*$
	be finite. Then $\constrsync{L \cup F} \leACZEROm \constrsync{L}$.	
\end{proposition}
\begin{proof}
	Looking into the proof of \cref{constr_sync:thm:finite},
	checking whether a word from a finite set
	synchronizes an input automaton can be incorporated into
	an $\leACZEROm$-reduction.
	Specifically, the reduction first checks whether
	an input automaton has a synchronizing word in~$F$
	and if it does, passes a \rYES-instance to $\constrsync{L}$,
	or otherwise passes the input to~$\constrsync{L}$.	
\end{proof}

\begin{corollary}
	For regular $L \subseteq \Sigma^*$
	and finite~$F \subseteq \Sigma^*$
	we have $\constrsync{L} \leACZEROm \constrsync{L \setminus F}$.
\end{corollary}
\begin{proof}
	This follows from \cref{constr_sync:prop:finite} since $L = ( L \setminus F ) \cup (L \cap F)$.	
\end{proof}